\section{Introduction}

**** Example Text Start ******
\gls{ai} has become an important field in computer science and software engineering. 
Recent progress in \gls{nlp} has been strongly influenced by the development of the 
\gls{transformer} architecture and the emergence of \gls{llm} systems.

Although \gls{llm} systems can achieve impressive performance in many language-related tasks, 
their internal behavior is often difficult to understand. This lack of transparency motivates 
the field of \gls{mechanisticinterpretability}, which aims to analyze the internal mechanisms 
of neural networks instead of treating them only as black-box models.

In this paper, \gls{sentimentanalysis} is used as a case study to investigate how a language 
model processes emotional information in text. For example, the model may classify a sentence 
such as ``The movie was surprisingly good'' as positive, while assigning a negative sentiment 
to a sentence such as ``The service was disappointing.''

The Transformer architecture introduced the self-attention mechanism, which became a foundation 
for many modern language models \cite{vaswani2017attention}. Later models such as GPT further 
demonstrated the effectiveness of large-scale pretraining for language understanding and generation 
\cite{radford2018improving}.

The aim of this work is to explore how internal elements such as \glspl{attentionhead}, 
\glspl{token}, and \glspl{activation} may contribute to sentiment-related behavior in a 
large language model.
**** Example Text END ******